\section{Proofs of the reduction lemmas}\label{app:proofs}

We use two facts. (F1) Put $L=\cl(\Nb(v))$. Then $\cl(G)=\cl(G-v)\cup(v*L)$ and the intersection is $L$. If $L$ is non-empty and contractible, both $L$ and $v*L$ are contractible; gluing along these simplicial subcomplexes shows that $\cl(G-v)\hookrightarrow\cl(G)$ is a homotopy equivalence. For an edge $uv$, put $L=\cl(\Nb(u,v))$. Then
\[
\cl(G)=\cl(G-uv)\cup((uv)*L),\qquad
\cl(G-uv)\cap((uv)*L)=(\partial(uv))*L.
\]
If $L$ is non-empty and contractible, both the join $(uv)*L$ and its attaching subcomplex $(\partial(uv))*L$ are contractible, and the same gluing argument applies. Thus deleting an $\IS$-deletable element preserves homotopy type and acyclicity. It also preserves membership in $\I$, directly by (I1) or (I3), since $\IS\subseteq\I$. The reverse additions preserve these properties as well. (F2) If $H\in\IS$ and $v$ (respectively $uv$) has neighbourhood (respectively common neighbourhood) in $\IS$ in $G$, with $G-v=H$ (respectively $G-uv=H$), then $G\in\IS$: prepend this deletion to a reduction of $H$. This implication does not assert that every legal deletion from a member of $\IS$ remains in $\IS$.

\begin{proof}[Proof of Lemma~\ref{lem:stuck}]
(a) Let $G$ be $\Ztwo$-acyclic, not in $\IS$, on $n_0$ vertices. If $G$ has an $\IS$-deletable vertex $v$, then $G-v$ is $\Ztwo$-acyclic (F1) and not in $\IS$ (F2), with $n_0-1$ vertices, contradicting the minimality of $n_0$. If $G$ has an $\IS$-deletable edge $uv$, then $G-uv$ is $\Ztwo$-acyclic and not in $\IS$ with the same number of vertices and one edge fewer, and $G=(G-uv)+uv$ is an (I4) addition. Induction on the number of edges shows that $G$ lies in the $\mathrm{I}_4$-closure of a graph with no $\IS$-deletable element, which is stuck; if $G$ has the minimum number of edges, it has no $\IS$-deletable edge and is itself stuck.

(b) Let $G\in\I\setminus\IS$ on $n_1$ vertices. If $v$ is $\IS$-deletable then $G-v\in\I$ (transformation (I1), since $\IS\subseteq\I$) and $G-v\notin\IS$ (F2), contradicting the minimality of $n_1$. If $uv$ is $\IS$-deletable then $G-uv\in\I\setminus\IS$ by (I3) and (F2), has one edge fewer and is $K_4$-free if $G$ is. Induction on the number of edges concludes as in (a); the graph reached is $\Ztwo$-acyclic because it is in $\I$.
\end{proof}

\begin{proof}[Proof of Lemma~\ref{lem:2conn}]
Let $G=A\cup_u B$ be stuck. Since $\cl(G)=\cl(A)\cup\cl(B)$ with $\cl(A)\cap\cl(B)=\{u\}$ contractible, the Mayer--Vietoris sequence gives $\widetilde H_*(G;\Ztwo)\cong\widetilde H_*(A;\Ztwo)\oplus\widetilde H_*(B;\Ztwo)$, so $A$ and $B$ are $\Ztwo$-acyclic. Because $u$ is a cut vertex, for every $w\in V(A)\setminus\{u\}$ we have $\Nb_G(w)=\Nb_A(w)$, and for every edge $e$ of $A$ we have $\Nb_G(e)=\Nb_A(e)$ (a common neighbour of the ends of $e$ lies in $A$). Since $G$ is stuck, no vertex $w\ne u$ and no edge of $A$ is $\IS$-deletable in $A$: $A$ is almost stuck at $u$, and likewise $B$. If $A$ had only two vertices, $A=uw$ and $\Nb_G(w)=\{u\}\cong K_1\in\IS$, contradicting that $G$ is stuck; so $|A|,|B|\ge 3$. If no almost stuck graph has at most $m$ vertices then $|A|,|B|\ge m+1$ and $|G|=|A|+|B|-1\ge 2m+1$, which proves the last statement.
\end{proof}

\begin{proof}[Proof of Lemma~\ref{lem:degrees}]
Let $uv$ be an edge in no triangle. Since $G$ is $2$-connected, $uv$ lies in a cycle $C$. The $1$-chain $C$ (over $\Ztwo$) is a cycle; the boundary of any $2$-chain has coefficient $0$ on $uv$ because no triangle contains $uv$, so $C$ is not a boundary and $\widetilde H_1(G;\Ztwo)\ne 0$, a contradiction. Hence every edge lies in a triangle and every $\Nb(u,v)$ is non-empty; since $G$ is stuck, $\Nb(u,v)\notin\IS$, so $|\Nb(u,v)|\ge 2$. Now fix $v$ and let $H=\Nb(v)$, a graph with $|H|\ge 1$ and $H\notin\IS$, in which every vertex $u$ satisfies $\Nb_H(u)=\Nb(u,v)\neq\emptyset$ with $|\Nb_H(u)|\ge 2$. A graph with at most three vertices in which every vertex has at least two neighbours is $K_3\in\IS$; hence $|H|\ge 4$, i.e.\ $\delta(G)\ge 4$. Finally, if $\deg v=n-1$ then $G$ is a cone with apex $v$, which is $0$-dismantlable and hence in $\IS$, a contradiction; so $\Delta(G)\le n-2$.
\end{proof}
